\chapter{Génese Dependente V: Abandonar o Desejo}

\lettrine{O}{\space surgimento} de \emph{dukkha} deve"-se ao apego aos desejos. E a percepção é que existe esta origem ou surgimento e que o desejo deve ser abandonado. Esta é a Segunda Nobre Verdade; é o conhecimento intuitivo do abandono.

Algumas pessoas pensam que tudo o que ensino é «aconteça o que acontecer, liberta"-te». Mas o ensinamento envolve uma investigação real do sofrimento: a percepção de libertar"-se ocorre através dessa compreensão. Portanto, «libertar"-se» não provém de um desejo de \emph{se livrar} do sofrimento -- isso não é «libertar"-se», pois não?

O \emph{vibhava"-tanha} -- ou desejo de nos livrarmos -- é bastante subtil, mas querer livrar"-nos das nossas impurezas é outro tipo de desejo. ``Deixar ir'' não é livrar"-se ou rejeitar com qualquer aversão. ``Deixar ir'' significa ser capaz de estar com o que é desagradável sem ficar preso à aversão -- porque a aversão é um apego. Se tiveres muita aversão, continuarás apegado. Medo, aversão -- tudo isto é agarrar"-se e apegar"-se.

A imparcialidade é a aceitação e a consciência das coisas tal como são, libertando"-se da aversão ao que é feio ou desagradável. Portanto, «deixar ir» não é uma frase engenhosa inventada como forma de descartar as coisas, mas sim uma visão profunda da natureza das coisas. «Deixar ir» é, portanto, ser capaz de suportar algo desagradável e não ficar preso à raiva e à aversão. A imparcialidade não é depressão.

Quantos de vocês ignoram e se recusam a reconhecer o caráter desagradável das funções dos vossos próprios corpos? Existem certas funções do corpo humano que não são bonitas, que na sociedade educada não mencionamos. Usamos todo o tipo de eufemismos e formas de nos desculparmos educadamente no momento apropriado, porque ninguém quer que a perceção de si mesmo esteja ligada a essas funções. Queremos que a nossa presença ou imagem esteja associada a algo agradável, interessante ou atraente. Queremos que a nossa fotografia seja tirada com flores num cenário atraente, não na casa de banho. Queremos disfarçar os processos naturais da vida, esconder as rugas, pintar o cabelo, fazer tudo para parecer mais jovens -- porque o envelhecimento não é atraente.

À medida que envelhecemos, perdemos o que é belo e atraente. A nossa reflexão deve ser realmente consciente da doença e da morte, do que é atraente e do que não é: a forma como as coisas são neste reino da consciência sensorial. Sendo uma entidade com órgãos sensoriais que entram em contacto com objetos -- que podem ser qualquer coisa, desde o mais belo e agradável até ao mais hediondo e feio -- experimentamos sensações. A sensação (\emph{vedana}) implica o dualismo do agradável, do doloroso e do neutro; isto aplica"-se a todos os sentidos -- paladar, tato, visão, audição, olfato e pensamento.

Utilizo essa palavra específica, «vedana», esse \emph{khandha}, como o conceito para toda essa atração/repulsa. Estamos a experimentar \emph{vedana}, estamos conscientes do agradável, do doloroso, do belo, do feio e do neutro através do corpo ou através do que ouvimos, cheiramos, provamos, tocamos ou pensamos. Até as memórias podem ser atraentes. Podemos ter memórias que são agradáveis, desagradáveis ou neutras. E se formos descuidados e agirmos a partir de \emph{avijja,} a visão do eu, a suposição inquestionável de que «eu sou» -- o atraente, o pouco atraente e o neutro são interpretados com desejo. Quero o belo, quero o agradável, quero ser feliz e bem"-sucedido. Quero ser elogiado, quero ser apreciado, quero ser amado. Não quero ser perseguido, infeliz, doente, menosprezado ou criticado. Não quero coisas feias à minha volta. Não quero olhar para o feio, estar rodeado do desagradável.

Consideremos as funções do nosso corpo. Todos sabemos que essas funções fazem simplesmente parte da natureza, mas não queremos pensar nelas como sendo \emph{minhas}. Tenho de urinar, mas ninguém gostaria de ficar na história como «Sumedho, o Urinador». «Sumedho, o Abade de Amaravati», isso já é aceitável. Quando escrever a minha autobiografia, ela estará repleta de coisas como o facto de eu ter sido discípulo de Ajahn Chah, de como eu era sensível quando era criança, inocente e puro -- talvez um pouco travesso de vez em quando, porque não quero ser visto como um boneco de pano. Mas, na maioria das biografias, as funções desagradáveis do corpo são simplesmente ignoradas. Não devemos andar por aí a pensar que devemos \emph{identificar"-nos} com estas funções, mas apenas começar a notar a tendência para não querer ser incomodados por elas, ou prestar atenção e observar muito daquilo que faz parte da nossa vida, a forma como as coisas são.

Na atenção plena, então, estamos a abrir a nossa mente a isto, à totalidade da vida, que inclui o belo, o feio, o agradável, o doloroso e o neutro. Assim, na nossa reflexão sobre o \emph{paticcasamuppada}, vemos que está ligado à Segunda Nobre Verdade.

É aqui que a sequência \emph{tanha"-upadana"-bhava} é mais útil como meio de investigar o apego. Apego, neste sentido, pode significar agarrar"-se por atração ou por aversão, tentando livrar"-se de algo. Apegar"-se com aversão é afastar; fugir é \emph{upadana}, assim como tentar agarrar o belo, possuí"-lo e mantê"-lo, buscando o desejável e tentando livrar"-se do indesejável.

Quanto mais contemplamos e investigamos \emph{o upadana}, mais surge a percepção: o desejo deve ser abandonado. Na Segunda Nobre Verdade, explica"-se que o sofrimento surge, deve ser abandonado e, então, através da prática do abandono e da compreensão do que realmente é o abandono, temos a terceira percepção da Nobre Verdade: o desejo foi abandonado: nós realmente sabemos o que é abandonar. Não é um abandono teórico, não é uma rejeição de nada, é a percepção real.

Ao discutir a Segunda Nobre Verdade, há as afirmações: «Existe a origem do sofrimento», «deve ser abandonado»; e a terceira compreensão: «foi abandonado». E é disso que se trata a prática -- cumprir essas três.

Isso aplica"-se a cada uma das Nobres Verdades: existe sofrimento; deve ser compreendido; foi compreendido. Estes são os três aspetos da compreensão da Primeira Nobre Verdade. A Segunda Nobre Verdade: existe a origem do sofrimento, que é o apego ao desejo; deve ser abandonado; foi abandonado. A Terceira Nobre Verdade: existe a cessação do sofrimento, \emph{nirodha}; deve ser realizada; e a terceira percepção: foi realizada. A Quarta Nobre Verdade: existe o Caminho Óctuplo, o caminho para sair do sofrimento; deve ser desenvolvido; foi desenvolvido. Isto é conhecimento intuitivo.

Quando pensas: «O que sabe um arahant?», é isto: ele sabe que existe sofrimento, sabe que o sofrimento deve ser compreendido, o arahant sabe quando o sofrimento foi compreendido; o arahant conhece a origem do sofrimento, sabe que deve ser abandonado, sabe que foi abandonado, etc. Estas são as doze percepções. É a isto que chamamos arahantidade -- o conhecimento de quem possui essas percepções.

\emph{Paticcasamuppada} é uma investigação muito aprofundada de todo o processo. Ora, o problema reside na apreensão dos cinco \emph{khandhas}. Os cinco \emph{khandhas} são dhammas\footnote{\emph{Dhamma:} significa «coisa» como parte do universo (ou seja, não pertencente a uma pessoa); Dhamma refere"-se ao ensinamento do Buda e à sua compreensão da Verdade Última.} -- devem ser estudados e investigados. São simplesmente como as coisas são. Não são o eu; são impermanentes. Saber isto, saber como as coisas são, é conhecer o Dhamma.

E assim, a apreensão do mundo condicionado como um eu baseia"-se na ilusão ou ignorância (\emph{avijja}): a ilusão de um eu como sendo os cinco \emph{khandhas}. E por causa disso vivemos as nossas vidas com base na ignorância. As atividades volitivas (\emph{sankhara}) decorrentes dessa ignorância interpretam tudo com o «eu sou» e a partir da apreensão dos desejos; o resultado é \emph{jara"-maranam} (envelhecimento e morte). Se eu apego o corpo como eu, então «eu» envelheço. O meu corpo tem 54 anos este ano, está flácido e enrugado. E a crença de que «estou a envelhecer» porque o corpo está a envelhecer é uma espécie de sofrimento, não é? Se não há sentido de eu, então não há sofrimento. Há uma apreciação do seu envelhecimento. Não há a sensação de que haja algo de errado com o corpo a envelhecer; é isso que ele deve fazer. Essa é a sua natureza. Não sou eu. Não é meu, e está a fazer o que deve fazer. Perfeito, não é? Ficaria chateado se ele começasse a ficar mais jovem! Imagine se eu começasse a ficar mais jovem, daqui a cinquenta anos estaria de volta às fraldas e teria de passar por tudo isso novamente.

O pensamento «Estou a envelhecer» não é triste. É uma forma convencional de falar sobre o corpo. Mas se é isto que eu penso que sou: «Eu sou o corpo, este é o meu corpo», então a ignorância condiciona \emph{o sankhara} e tudo provém desse «Estou a envelhecer, quero ser jovem, quero viver uma vida longa, não me chames de velho, seu jovem insolente!» Porquê? Por causa da identificação com o corpo.

Também reflito que vou morrer. «Isso é algo mórbido, nem vamos falar de morte. Claro que todos vamos morrer, mas isso está muito longe.» Quando se é jovem, pensa"-se na morte como algo tão distante -- «vamos aproveitar a vida.» Mas quando alguém que conhecemos morre, ou quando estamos prestes a morrer, a morte pode ser muito assustadora. E tudo isso vem do apego à identidade com este corpo.

Depois, claro, há todas as visões, sentimentos, memórias e preconceitos que temos \emph{(vedana, sanna, sankhara).} Não só sofremos com a identificação com o corpo, mas também quando nos apegamos ao belo e aos sentimentos: «Só quero o belo, só quero o agradável; não quero ver o feio; quero ter música bonita e nenhum som feio, apenas os cheiros perfumados\ldots{}»

Apegamo"-nos à forma como o mundo \emph{deveria} ser; opiniões sobre a Grã"-Bretanha, a França e os EUA. Os apegos a estas visões, opiniões e perceções constituem a sequência \emph{vedana"-sanna"-sankhara} dos cinco \emph{khandhas}. E podemos apegar"-nos a tudo isso em termos de «eu». «É a minha visão, o que penso, o que quero e não quero, o que deveria ser e o que não deveria ser.» Assim, sentimos sofrimento, angústia, desespero, depressão, tristeza e lamentação devido a essa ilusão do eu.

A compreensão da Segunda Nobre Verdade é que existe uma origem para este sofrimento. Não é permanente. Não é absolutamente certo que algo surja. O surgimento do \emph{dukkha} deve"-se ao apego do desejo. É possível perceber o desejo porque é \emph{dhamma}: surge e cessa.

Pode"-se ver o desejo que surge para buscar o belo e o agradável no plano sensual através dos olhos, ouvidos, nariz, língua, corpo e mente. \emph{Kama"-tanha} é o desejo sensual. O desejo sensual busca sempre algum tipo de experiência prazerosa ou, pelo menos, emocionante. \emph{Kama"-tanha}: pode"-se vê"-lo no movimento que temos de nos dirigir para os prazeres sensoriais e, em seguida, agarrá"-los.

\emph{Bhava"-tanha} é o desejo de se tornar. Isto tem a ver com querer tornar"-se algo. Em última análise, uma vez que não sabemos quem somos, o nosso desejo é alcançar, realizar e \emph{tornar"-se} algo. Nesta Vida Santa, o \emph{bhava"-tanha} pode ser muito forte. Sente que está aqui para \emph{se tornar} iluminado e alcançar e atingir algo. Tudo isso soa muito bem. Mas até mesmo o desejo de se tornar iluminado pode vir deste \emph{avijja}, desta visão do eu. «Vou me tornar iluminado. Vou me tornar o primeiro arahant americano; estou farto deste mundo, quero me tornar iluminado para não ter que renascer novamente. Não quero passar pela infância outra vez. Não quero nada disso. Quero tornar"-me algo em que já não seja preciso nascer.» Isso pode ser \emph{bhava"-tanha/vibhava"-tanha} -- andam de mãos dadas. Para se tornar algo, tem de se livrar das coisas de que não gosta e que não quer. «Vou livrar"-me das minhas impurezas e quero livrar"-me dos meus maus hábitos e livrar"-me dos meus desejos» -- e tudo isto também soa muito correto. As impurezas são más -- livra"-te delas.

Assim, na Vida Santa, há muito \emph{vibhava"-tanha}. Podemos viver esta vida apenas para nos livrarmos de coisas e para nos tornarmos algo ao livrarmo"-nos de algo. Reparem, então, que a Segunda Nobre Verdade é a compreensão de que o desejo deve ser abandonado, deve ser posto de lado. Não é uma rejeição do desejo, mas uma compreensão; você o deixa ir. Porque, caso contrário, é \emph{vibhava"-tanha,} o desejo de se livrar do desejo. Reconheça-o, veja-o, mas não tire conclusões a partir dele. Se você está vindo da ignorância, o seu desejo diz: «Eu quero tornar"-me um ser iluminado e não devo pensar assim, não devo ter o desejo de me tornar um Buda; não devo querer tornar"-me nada.» Tudo isso pode ser fruto da ignorância condicionando as formações mentais \emph{(avijjapaccaya sankhara)}. Então, surge o conhecimento intuitivo: «O desejo deve ser abandonado.»

Tudo isto soa muito certo, de certa forma, quando dizemos «não devemos apegar"-nos a nada», mas isso também pode provir de \emph{avijjapaccaya sankhara}. «Não devo apegar"-me a nada» é, em grande medida, uma afirmação de mim mesmo como alguém que está apegado a algo e não deveria ser assim. Eu deveria ser diferente. Portanto, isso é apenas uma armadilha da mente, não é um verdadeiro insight sobre \emph{kama"-tanha}, \emph{bhava"-tanha, vibhava"-tanha}. Refletam sobre o que é o apego. Se estiverem realmente apenas a deitar coisas fora, essa não é a forma de resolver este problema, pois não? Tens de examinar verdadeiramente \emph{kama"-tanha, bhava"-tanha, vibhava"-tanha,} para que não tenhas uma compreensão sobre o desapego. Limitar"-te-ás a assumir uma posição contra o apego, o que é outro tipo de apego.

Por isso, examina, analisa o apego. Isto funciona de uma forma muito mais subtil e realista do que simplesmente formar uma opinião de que não deves estar apegado a nada.

Lembro"-me de um psiquiatra que vivia em Banguecoque, que costumava tirar o relógio de pulso de alguém\ldots{} e a pessoa ficava chateada e ele dizia: «Estás apegado ao teu relógio de pulso.» Depois, tirava o seu próprio relógio e atirava-o fora para provar que não estava apegado. Ele gabava"-se disso comigo. Eu disse: «Não percebeste o essencial. Estás apegado à ideia de que não estás apegado ao teu relógio de pulso.» Isso não é libertar"-se, pois não? Deitar isto fora como um espertalhão e dizer: «Tu estás apegado, eu não, eu deitei o meu fora.» Há muito de ego nisso. Olha para mim, não estou apegado a estas coisas materiais miseráveis.» Podes sentir"-te bastante orgulhoso de não estares apegado. Com reflexão, vemos o apego e não precisamos de nos livrar das coisas, mas podemos não estar apegados a elas, podemos libertar"-nos delas -- não atirando"-as fora, mas compreendendo o sofrimento que advém do apego.

À medida que compreendes a paz do desapego, do deixar ir, a Segunda Nobre Verdade conduz à Terceira. Quando deixas ir algo, estás ciente de que não há apego aos cinco \emph{khandhas}. Há a consciência de que o desejo foi largado. Então surge a compreensão da Terceira Nobre Verdade da cessação. Há cessação. Esta cessação deve ser compreendida.

À medida que percebemos cada vez mais a cessação, começamos a notar o desapego. Poucos de vós estão conscientes do desapego. Normalmente, estais conscientes por estardes apegados às coisas. Um ser humano totalmente iludido só se sente vivo através do apego e do desejo. Contemplai que, quando não estais presos ao apego aos cinco \emph{khandhas}, não vos sentis vivos, não sois ninguém. Ter problemas neuróticos faz com que as pessoas se sintam interessantes e vivas. «Tenho neuroses fascinantes decorrentes de todo o tipo de traumas na primeira infância.» Portanto, não é o Sumedho, o Urinador, é o Sumedho, o Neurótico Interessante, o Místico ou o Sumedho, o Abade -- estas são condições às quais podemos estar apegados. Perceber a cessação permite"-vos deixar o eu cessar; há um desapego. A compreensão do desapego é a cessação, que tudo o que surge cessa. E a cessação é observada. A cessação deve ser compreendida.

Portanto, a nossa prática é a \emph{de perceber} a cessação. É disso que falamos quando falamos de vazio: percebemos a mente vazia onde não há eu. Não há a sensação de que a mente seja alguém. Assim que pensas nisto como «a minha» mente, se te agarrares a esse pensamento, então estás novamente iludido. Mas mesmo que tenhas «a minha» mente e a vejas como aquilo que surge e cessa sem te agarrares a ela, então é apenas uma condição. Não há sofrimento nisso, é pacífico.

Quando não há eu, há paz. Quando há «eu» e «meu», então não há paz. Preocupação, ansiedade, o que são? São todas provenientes do «eu» e do «meu». Quando se deixa ir, então há cessação do «eu» e do «meu». Há paz, calma, clareza, desapego, vazio.

Observo que, quando não há ego, nem apego, a forma de nos relacionarmos com os outros é através de \emph{metta} (bondade), \emph{karuna} (compaixão), \emph{mudita} (alegria solidária) e \emph{upekkha} (serenidade). Estas não provêm do ego nem \emph{da avijja.} Não é que exista a ideia de que «tenho de ter mais metta por todos porque tenho muita aversão e não devia. Devo ter bondade amorosa por todos os seres. Devo sentir compaixão. Às vezes, só me apetece matar toda a gente. Devo sentir muita \emph{metta, mudita,} ser gentil, alegre e solidário com as pessoas. Devo ser sereno também.» Os brahma"-viharas\footnote{\emph{Brahma"-viharas:} As «moradas divinas» da bondade, compaixão, alegria solidária e serenidade.}, como ideias para uma pessoa egoísta, não são a prática real. O desejo de se tornar alguém que tem muito \emph{metta} e \emph{karuna} e todo esse tipo de coisa ainda é \emph{bhava tanha.}

Mas à medida que as ilusões do eu vão desaparecendo, essa passa a ser a forma natural de nos relacionarmos. Não nos tornamos um zombie vazio por compreender o Dhamma. Continuamos a relacionar"-nos uns com os outros, mas é através da bondade, da compaixão, da alegria solidária e da serenidade, em vez de através da ganância, do ódio e da ilusão.

O que é que os seres humanos altruístas, geralmente, manifestam na sociedade? Pode explicar"-se \emph{metta, mudita, karuna e upekkha} como qualidades que se manifestam através de seres humanos altruístas. Depois, aplique isso à nossa própria prática agora. Quando há \emph{vijja} -- saber e ver claramente --, isso dá total oportunidade para a prática da bondade, compaixão e o resto. Mas não sou eu, não é meu; não é Sumedho, o Ajahn cheio de metta, Sumedho, o Bom Rapaz, mas sim Sumedho, o Urinador. Assim que as ilusões de Sumedho se afastam e cessam, a bondade, a compaixão, a alegria solidária e a serenidade podem manifestar"-se. É por isso que o estado humano é uma grande bênção: quando a visão do eu é abandonada, o que resta é uma grande bênção.

Mas não sou eu. «Eu não sou uma grande bênção. Tudo o que posso fazer neste eu convencional é libertar"-me da ilusão. Estar atento e não me apegar às coisas, ver com clareza -- é isso que posso fazer. Essa é a prática das Quatro Nobres Verdades e o desenvolvimento do Caminho Óctuplo. Resumindo"-se a essa visão vigilante e consciente das coisas com clareza. O que acontece depois depende de outras coisas. Já não há necessidade de andar por aí a tentar tornar"-me «Sumedho, o Bom Rapaz». Esta bondade pode manifestar"-se através desta forma se não houver ilusão; e isso não é de todo uma conquista pessoal, é apenas a forma como as coisas são. A forma como acontece. É o Dhamma.

\photoPage{image-014-sweeping-leaves.jpg}{Bhikkhu a varrer folhas, Wat Pah Nanachat, Tailândia}{Foto de Ping Amranand}


